\section{Numerical facts and the chapter protocol}
\label{app:lvnumerics}

\subsection{The protocol}
\label{app:lvprotocol}

Every fit in this chapter --- synthetic and real --- uses the following
recipe; no constant is tuned per name.

\emph{Vertex grid.}  Maturity rows at $\tau=0$ and at every quoted expiry.
Strike columns delta-spaced: the ladder of put deltas
$\{1,2,5,10,25,40\}\%$, the same ladder mirrored on the call side, and
the money, mapped to
strikes through
$y=\exp\big(\sigma_A\sqrt{\tau_{\max}}\;\PhiN^{-1}(\cdot)\big)$ with
$\sigma_A$ the longest expiry's ATM volatility; the two end columns are
extended to cover the pooled quoted range, so no quote falls outside the
hull.  Coverage rule: every expiry must own at least $6$ columns inside
its own traded range (\cref{sec:lvinfo}), enforced as follows: list
the columns falling inside the expiry's range together with the
range's two endpoints, and while the widest gap between consecutive
entries exceeds one fifth of the range, split that gap with a new
column; this is what densifies the short-dated money in
\cref{fig:lvsheet}.

\emph{Lattice.}  Strike step
\[
\Delta y=\min\big(0.01,\;0.15\,\min_e\sigma_e\sqrt{\tau_e}\big),
\]
capped at $800$ nodes over $[0,\,\max(2.5,\,1.05\,y_{n_y})]$ with
$y_{n_y}$ the highest vertex column, and $y=1$ a node; every quoted
expiry is a time node and each maturity interval receives
$\max(8,\lceil\Delta\tau_{\rm int}/0.01\rceil)$ implicit steps.
The refined (audit) operator halves the strike step and quadruples each
interval's step count.

\emph{Box and objective.}  Adaptive box
\[
v_{\rm lo}=\min\big(0.0025,\,(\tfrac12\min_e\sigma_{{\rm ATM},e})^{2}\big),
\qquad
v_{\rm hi}=\min\big(\max(0.36,\,(3\,\sigma_{\max})^{2}),\,16\big)
\]
--- the floor a fraction of the smallest ATM implied per the averaging
fact, the cap a multiple of the largest implied in view.  Seed and reference
$\theta^{0}=\theta_{\rm ref}$ flat at the median ATM variance.  Weights
$\lambda=10^{-2}$, $\lambda_0=1$; residual scale $\eta_q$ equal to one
volatility percent of the quote's Black vega, floored at $10^{-3}$.  The
snapshot carries no variance-swap quotes, so the wing multiple stays at
its default $a=0$ (flat clamp) throughout.

\emph{Solver.}  Trust-region reflective least squares
\citep{BranchColemanLi1999} with the analytic Jacobian of
\cref{prop:lvtangent}, termination tolerances $10^{-8}$, at most $200$
objective evaluations; the fits of \cref{sec:lvexamples} converged in
\MacLvSpyEvals\ (SPY), \MacLvNvdaEvals\ (NVDA) and \MacLvRtEvals\
(synthetic) evaluations from their flat seeds.

\subsection{The synthetic market}
\label{app:lvsynthetic}

The running example's truth is the smooth skewed field
\[
\sigma_{\rm loc}(\tau,y)
=0.21-0.10\,\tanh\!\Big(\frac{y-1}{0.22}\Big)
+0.03\,e^{-2\tau},
\qquad v=\sigma_{\rm loc}^{2},
\]
priced by an \emph{independently implemented} dense march (uniform
$\Delta y=0.0025$ on $[0,3]$, $\Delta\tau=0.002$, fully implicit) --- a
second discretization on different grids, so the round trip of
\cref{fig:lvrecovery} cross-checks the chapter's march rather than
validating it against itself.  Quotes are read at the four expiries
$\tau\in\{0.10,0.25,0.50,1.00\}$, at $31$ strikes per expiry uniform in
the standardized moneyness $\zeta\in[-2.2,2.2]$ of \cref{sec:mcs}, with
$\sigma_{\rm ref}=0.24$, i.e.\ $k=0.24\sqrt{\tau}\,\zeta$ (the coarse arm
of \cref{fig:lvwrongway}a uses $13$).  The ``noise'' of
\cref{sec:lvillposed} is a deterministic alternating
$\pm50$\,bp volatility ripple --- no randomness anywhere in the chapter's
figures.  The synthetic sheet has rows $\{0\}\cup\{\text{expiries}\}$ and
the delta-spaced columns at scale $0.24$ plus one deliberate deep-put
column at $y=\MacLvIdentColY$ for the experiment of
\cref{fig:lvidentify}; its box is $[0.005,0.36]$.

\subsection{The adjoint}
\label{app:lvadjoint}

The tangent route of \cref{prop:lvtangent} prices the Jacobian at
$O(N_\tau N_y\,n_\theta)$ per evaluation, for a march of $N_\tau$ steps on
$N_y$ lattice nodes.  The classical alternative makes the cost of a
\emph{gradient} independent of $n_\theta$; deriving it costs a few lines,
so the chapter records it for the grids on which it will one day pay.
The derivation follows the standard optimal-control pattern ---
adjoin the marching constraints with multiplier vectors $p^{n}$, the
\emph{adjoint states}, require stationarity, and read the result as a
backward march; a reader meeting adjoints for the first time can treat
\eqref{eq:lvadjoint}--\eqref{eq:lvadjgrad} as a recipe and check it
numerically against the tangent Jacobian.

\begin{proposition}[Adjoint: one backward sweep]
\label{prop:lvadjoint}
Let $\mathcal J(\theta)$ be a scalar objective depending on the march
through the observed values $\mathcal O_nU^{n}$, with $\mathcal O_n$ fixed
observation rows.  Define adjoint states by the backward recursion
\begin{equation}
(\mathcal M^{n})^{\!\top}p^{n}
=\mathcal O_n^{\top}\,
\frac{\partial\mathcal J}{\partial(\mathcal O_nU^{n})}+p^{\,n+1},
\qquad p^{\,N+1}=0 .
\label{eq:lvadjoint}
\end{equation}
Then
\begin{equation}
\frac{\partial\mathcal J}{\partial\theta_\ell}
=\sum_{n}\Delta\tau\,(p^{n})^{\!\top}\,
\mathcal L[\mathcal H_\ell]\,U^{n},
\label{eq:lvadjgrad}
\end{equation}
with the boundary-inclusive product of \cref{prop:lvtangent}, at the cost
of one backward solve, $O(N_\tau N_y)$, plus the $O(n_\theta)$
accumulation of \eqref{eq:lvadjgrad}.
\end{proposition}

\begin{proof}
Form the Lagrangian (fraktur $\mathfrak L$, kept visually apart from
the stencil operator $\mathcal L$)
$\mathfrak{L}=\mathcal J-\sum_n (p^{n})^{\top}
\big(\mathcal M^{n}U^{n}-U^{n-1}-b^{n}\big)$, with $b^{n}$ the boundary
terms of \cref{prop:lvtangent}.  Stationarity in each $U^{n}$ --- which
appears in step $n$'s system and in step $n{+}1$'s right-hand side ---
gives \eqref{eq:lvadjoint}; the $\theta$-derivative of $\mathfrak{L}$ then
retains the explicit dependence of $\mathcal M^{n}$ and $b^{n}$, which
combine as in \cref{prop:lvtangent} into the boundary-inclusive
\eqref{eq:lvadjgrad}.  Structurally this is summation by parts: the
transpose of a lower block-bidiagonal system is upper block-bidiagonal,
i.e.\ a backward march with transposed tridiagonal factors.
\end{proof}

\subsection{Measured cleanliness}
\label{app:lvdiag}

The diagnostics quoted in \cref{sec:lvexamples} are computed on the
fitting operator's price grid: the butterfly proxy is the minimum divided
second difference $\big(U_{i+1}-U_i\big)/\Delta y_i^{+}
-\big(U_i-U_{i-1}\big)/\Delta y_i^{-}$ over all interior nodes and all
expiries; the calendar proxy is the minimum of
$U^{(e+1)}-U^{(e)}$ over adjacent expiries at fixed $y$ (the lattice
prices are already in the normalized units of \cref{sec:calendar}, so the
fixed-$y$ comparison is the right one); and the bounds check verifies
$0\le U\le1$ to the same tolerance.  For both calibrated sheets all three
sit at solver rounding (\cref{sec:lvexamples}); they are measured per fit
precisely because \cref{rem:lvscope} declines to promote them to a
theorem.
